\chapter{Radiation, skin temperature and surface exchanges}
\label{cap:superficie}

Over the verification sample the archive carries more downward solar radiation at the surface than
ERA5 (\num{1.136} times), less total cloud cover (\num{0.783} times), and a skin temperature over
land \SI{1.64}{K} colder. Radiation, clouds and the surface energy budget must therefore be read
together.

\section{Radiation}
\label{sec:radiazione}

The 13 radiative fluxes of \cref{tab:radiazione} are fields accumulated over the hour ending at the validity time, in
\si{J/m^2} (\cref{sec:tempo}).

\begin{table}[htbp]
  \centering
  \small
  \caption{Radiative fluxes, accumulated over the hour ending at the validity time.}
  \label{tab:radiazione}
  \begin{tabular}{@{}lrl@{}}
    \toprule
    Field & paramId & Description \\
    \midrule
    \campo{ssrd}  & 169    & surface downward solar radiation \\
    \campo{ssrdc} & 228129 & surface downward solar radiation, clear sky \\
    \campo{ssr}   & 176    & surface net solar radiation \\
    \campo{ssrc}  & 210    & surface net solar radiation, clear sky \\
    \campo{strd}  & 175    & surface downward thermal radiation \\
    \campo{strdc} & 228130 & surface downward thermal radiation, clear sky \\
    \campo{str}   & 177    & surface net thermal radiation \\
    \campo{strc}  & 211    & surface net thermal radiation, clear sky \\
    \campo{tisr}  & 212    & incident solar radiation at the top of the atmosphere \\
    \campo{tsr}   & 178    & net solar radiation at the top of the atmosphere \\
    \campo{tsrc}  & 208    & net solar radiation at the top of the atmosphere, clear sky \\
    \campo{ttr}   & 179    & net thermal radiation at the top of the atmosphere \\
    \campo{ttrc}  & 209    & net thermal radiation at the top of the atmosphere, clear sky \\
    \bottomrule
  \end{tabular}
\end{table}

Solar radiation at the surface is higher than in ERA5: the downward flux by \SI{8.5}{\percent} under
clear sky and by \SI{13.6}{\percent} under all-sky conditions (\cref{tab:rad-era5}) \parencite{issue45}. The
simulation has no aerosol (\campo{aer_opt=0}), whereas ERA5 uses the aerosol climatology of
\textcite{tegen1997}. The clear-sky flux isolates the cloud-free atmosphere, and the thermal fields,
on which aerosol does not act, agree with ERA5 within \SI{1.1}{\percent}; the comparison attributes
the clear-sky excess to the absence of aerosol. With clouds the excess grows, because the total
cloud cover is \num{0.783} times that of ERA5 (\cref{sec:colonne}) and lets more radiation through.
For solar energy applications, or for training against pyranometer observations, the archive
carries a systematic excess of this size, which is a property of the simulation and which the
conversion cannot remove.

\begin{table}[htbp]
  \centering
  \small
  \caption{Ratio of means between CHAPTER and ERA5 over the verification sample, with the CHAPTER
  fields aggregated onto the ERA5 grid.}
  \label{tab:rad-era5}
  \begin{tabular}{@{}lS@{}}
    \toprule
    Field & {CHAPTER / ERA5} \\
    \midrule
    \campo{ssrdc} & 1.085 \\
    \campo{ssrd}  & 1.136 \\
    \campo{ssrc}  & 1.116 \\
    \campo{ssr}   & 1.178 \\
    \campo{tsrc}  & 1.031 \\
    \campo{strd}  & 0.989 \\
    \campo{strdc} & 0.996 \\
    \campo{ttr}   & 1.009 \\
    \campo{ttrc}  & 0.990 \\
    \bottomrule
  \end{tabular}
\end{table}

\section{Skin temperature}
\label{sec:pelle}

The simulation did not write the model skin temperature (\campo{TSK}), and the archive
reconstructs it as \campo{skt} from the outgoing thermal radiation at the surface. The upward flux
$F^{\uparrow}$ (\campo{LWUPB} in WRF) is
\begin{equation}
  F^{\uparrow} = \varepsilon \sigma T_s^4 + (1 - \varepsilon) F^{\downarrow},
  \label{eq:lwup}
\end{equation}
where $\varepsilon$ is the surface emissivity, $\sigma$ the Stefan--Boltzmann constant, $T_s$ the
skin temperature and $F^{\downarrow}$ the downward thermal flux at the surface (\campo{GLW}). The
emissivity is derived from the simulation itself. RRTMG also computes the clear-sky fluxes, with
the same emissivity and the same skin temperature and a different downward flux. Subtracting the
two forms of \cref{eq:lwup} eliminates the temperature:
\begin{equation}
  \varepsilon = 1 - \frac{F^{\uparrow} - F^{\uparrow}_\text{cs}}{F^{\downarrow} - F^{\downarrow}_\text{cs}},
  \label{eq:emissivita}
\end{equation}
where the subscript cs denotes clear sky. The expression is an identity, with no reference
temperature and no estimated parameters. Measured in this way, the category emissivity is the same
in every cell with the same dominant category, with an interquartile range of zero to the fourth
decimal place, and equals \num{0.98000} over water. The values per category coincide with the
\campo{LEMI} column of the \campo{MODI-RUC} section of \campo{VEGPARM.TBL} \parencite{wrf411}, from
\num{0.850} on barren soil to \num{0.980} over water.

With snow the emissivity switches to \num{0.98} for all categories. The archive applies \num{0.98}
where the snow cover \campo{SNOWC} exceeds the snow-cover threshold, \num{0.015}
(\SI{1.5}{\percent}), and below the threshold blends the two values. The snow-cover threshold is
the only parameter of the reconstruction not read from the simulation: it is the value that minimizes
the skin-temperature error in a sweep over the 34 validity times of the verification sample.

The accuracy depends on the season (\cref{tab:skt}). Over open water, where the model skin
temperature is the sea surface temperature, the reconstruction recovers it with a root-mean-square error of
\SI{0.0005}{K} on about \num{910000} points, at every validity time of the verification sample.
Over land the reconstructed emissivity coincides with the model one within $10^{-4}$ on
\SI{99.98}{\percent} of the points in July and on \SI{93.05}{\percent} in January. The whole of the
difference between the seasons comes from the snow rule. Over snow-free land the reconstructed skin
temperature differs from the temperature of the soil node at \SI{0}{cm} (\campo{TSLB}) by a
root-mean-square difference of \SI{0.026}{K}, on 4 September 2024 at 16 UTC.

\begin{table}[htbp]
  \centering
  \small
  \caption{Accuracy of the reconstructed skin temperature against the model state.}
  \label{tab:skt}
  \begin{tabularx}{\textwidth}{@{}Xl@{}}
    \toprule
    Quantity & Value \\
    \midrule
    open water, root-mean-square error (maximum absolute error) & \SI{0.0005}{K} (\SI{0.002}{K}) \\
    land, points with emissivity within $10^{-4}$ & \SI{99.98}{\percent} in July, \SI{93.05}{\percent} in January \\
    land, root-mean-square error in July & \SI{0.0012}{K} \\
    land, root-mean-square error in January and February & \SI{0.084}{K} to \SI{0.119}{K} \\
    land, January and February, 99th-percentile absolute error (maximum) & \SI{0.81}{K} (\SI{2.79}{K}) \\
    \bottomrule
  \end{tabularx}
\end{table}

Against ERA5 the skin temperature over land is \SI{1.64}{K} colder on average over the 34 validity
times of the verification sample. The difference, CHAPTER minus ERA5, ranges from \SI{-0.5}{K} to
\SI{-1.2}{K} among the 00 UTC validity times\todo{00 UTC validity times to be re-measured on the files of the previous day's run (issue 72)} and from \SI{-0.8}{K} to \SI{-3.24}{K} among the
12 UTC ones, with the coldest value on 15 September 2024. Over the sea the difference is
\SI{+0.01}{K}, and the correlation with ERA5 is between \num{0.95} and \num{0.98} over both land and
sea \parencite{issue57}. Over the sea, where CHAPTER and ERA5 use almost the same sea surface temperature,
the difference of \SI{0.01}{K} rules out the reconstruction, the emissivity and the encoding as
causes. The error of the reconstruction over land, at most \SI{0.119}{K} in root-mean-square error,
is more than ten times smaller than the mean difference from ERA5.

The difference over land has the opposite sign to the one that radiation and clouds suggest.
More solar radiation at the ground and fewer clouds bring energy to the surface, yet the surface is
colder, above all at 12 UTC, when the excess of solar radiation is largest. The partitioning of
energy at the surface therefore remains unexplained. Explaining it requires the surface energy
budget, whose three closing terms, the sensible heat flux, the latent heat flux and the ground heat
flux, are not in the archive (\cref{cap:limiti}). Without these fields four possible causes cannot
be separated: the surface scheme may send more energy into latent heat, the soil-moisture climate
may differ from that of ERA5, the orography resolved at \SI{3}{km} raises the mean elevation of a
cell aggregated onto the ERA5 grid, or the difference is a model error. Users who compare
\campo{skt} with satellite surface temperatures, or correct its bias, must use the measured
difference and not an estimate derived from radiation and clouds.

\section{Roughness length, friction velocity and stress}
\label{sec:rugosita}

The roughness length \campo{fsr} is WRF's local roughness length \campo{ZNT}: a table value per
dominant category, replaced under snow. It is \SI{0.70}{m} for forests, \SI{0.15}{m} for savannas
and grasslands, \SI{0.065}{m} for barren soil, \SI{1.00}{m} on the \num{19867} urban points, and
drops to \SI{0.011}{m} where snow covers the surface. The ERA5 roughness length is instead an
effective roughness length, which contains the orographic drag that the IFS parameterizes at
\SI{31}{km}; at \SI{3}{km} the orography is resolved. Over the sea both roughness lengths follow
the Charnock formula. Over land they differ by a factor of two to three (\cref{tab:rugosita}), and a
drag coefficient computed from the archive's \campo{fsr} differs from one computed from the ERA5
roughness length through the definition of the roughness length alone \parencite{issue56}.

\begin{table}[htbp]
  \centering
  \small
  \caption{Roughness length and friction velocity: range of the means per validity time over six
  validity times spread across the seasons, and ratios to ERA5.}
  \label{tab:rugosita}
  \begin{tabularx}{\textwidth}{@{}lXXXX@{}}
    \toprule
    Field & CHAPTER & ERA5 & Ratio of means & Ratio of medians \\
    \midrule
    \campo{fsr}, land  & \SIrange{0.188}{0.217}{m}   & \SIrange{0.375}{0.394}{m}   & \numrange{0.48}{0.55} & \numrange{0.32}{0.48} \\
    \campo{fsr}, sea   & \SIrange{0.0089}{0.0091}{m} & \SIrange{0.0103}{0.0105}{m} & \numrange{0.85}{0.88} & \numrange{1.5}{1.9} \\
    \campo{zust}, land & \SIrange{0.44}{0.52}{m/s}   & \SIrange{0.23}{0.28}{m/s}   & \numrange{1.62}{2.06} & \\
    \campo{zust}, sea  & \SIrange{0.21}{0.34}{m/s}   & \SIrange{0.19}{0.31}{m/s}   & \numrange{1.08}{1.12} & \\
    \bottomrule
  \end{tabularx}
\end{table}

The friction velocity \campo{zust} is the model's own. The ratio of means to ERA5 is between
\num{1.62} and \num{2.06} over land and between \num{1.08} and \num{1.12} over the sea, where the
correlation is between \num{0.85} and \num{0.90}. A roughness length between \num{0.32} and
\num{0.55} times that of ERA5, together with a friction velocity between \num{1.62} and \num{2.06}
times that of ERA5, rules out a simple scale error on the same roughness length. At \SI{3}{km} the
terrain-forced wind is resolved, and the model friction velocity follows it.

The surface stress components \campo{iews} and \campo{inss}, in \si{N/m^2}, are the magnitude
$\rho u_*^2$, with $\rho$ the air density and $u_*$ the model friction velocity, projected onto the
direction of the \SI{10}{m} wind. Alignment with the \SI{10}{m} wind is the only assumption. With
$\tau$ the vector of components \campo{iews} and \campo{inss}, the identities $|\tau| = \rho u_*^2$
and $u_* = \sqrt{|\tau|/\rho}$ hold within \SI{1.8e-6}{N/m^2} and within \SI{2.7e-4}{m/s} at every
point of six validity times.
